\documentclass[10pt,a4paper]{amsart}
\usepackage[margin=19mm]{geometry}
\usepackage{amsmath,amssymb,mathtools}
\usepackage[scr=rsfs,cal=euler]{mathalfa}
\usepackage{xfrac}
\usepackage[unicode,hidelinks,bookmarks=false]{hyperref}
\newcommand{\ie}{i.e.\ }
\newcommand{\cf}{cf.\ }
\newcommand{\resp}{respectively\ }
\newcommand{\Subsection}{\subsection}
\providecommand{\nobreakdash}{\nobreak\hbox{-}}
\setlength{\emergencystretch}{2em}
\multlinegap0pt
\usepackage{bm}
\usepackage{bbm}
\usepackage{dsfont}
\usepackage{xspace}
\usepackage{cancel}
\usepackage{mathtools}
\usepackage[shortlabels]{enumitem}
\usepackage{amsthm}
\makeatletter
\@ifundefined{theorem}{\newtheorem{theorem}{Theorem}}{}
\@ifundefined{lemma}{\newtheorem{lemma}[theorem]{Lemma}}{}
\makeatother
\def\gL{\mathfrak L}
\def\ga{\mathfrak a}
\title[A $p$-adic adjoint $L$-function and the ramification locus o]{A $p$-adic adjoint $L$-function and the ramification locus of the Hilbert modular eigenvariety}
\author{Baskar Balasubramanyam, John Bergdall, Matteo Longo}
\date{Source: arXiv:2011.05237v3 (CC BY 4.0)}
\begin{document}
\maketitle

\noindent\emph{Source and licence.} A $p$-adic adjoint $L$-function and the ramification locus of the Hilbert modular eigenvariety. Version arXiv:2011.05237v3 (CC BY 4.0), distributed under the Creative Commons Attribution 4.0 International licence, as recorded in the arXiv licence metadata for this exact version and shown on the abstract page https://arxiv.org/abs/2011.05237v3. Full rights notice: licenses/arxiv-2011.05237-CC-BY-4.0.txt.\par\smallskip

\noindent\textit{Editorial scope.} Counted unit: the Lemma whose source label is \texttt{lem:base-change-lemma} in the article (source0.tex lines 2311--2322, source label \texttt{lem:base-change-lemma}), delivered together with the article's own complete original proof of it (source0.tex lines 2323--2362, one \texttt{proof} environment of 40 lines). No mathematics is written, repaired or re-derived: statement and proof are the article's own text line for line. Required context retained uncounted: none. Every definition, display and label of those lines is retained unchanged. Omitted from this package and not redistributed: 1--2310, 2363--2899. Nothing inside a retained excerpt is omitted or altered: every retained body is delivered exactly as the article's lines are written, so no omission receipt of any kind accompanies this package. Citations: the retained excerpts carry no citation command at all, so no bibliography entry of the article is transcribed and no reference is redistributed. Reference adaptation: none was needed and none was made; every reference inside the delivered unit resolves inside it, so no unresolved reference or citation is produced. Printed slips kept literal: no printed word, symbol, hypothesis or number of the counted statement or of its proof was altered. Layout and class: the article's own class is \texttt{amsart} with packages including inputenc, amsthm, amsmath, amssymb, oldgerm, mathrsfs, geometry, verbatim, color, bbm, xy, stmaryrd, lmodern, graphicx; the wrapper is the lane's pinned \texttt{amsart} editorial wrapper at 10pt on A4, which restates the theorem-like environments the retained text needs and adds the article's own macro definitions that the retained text needs (\texttt{\textbackslash gL}, \texttt{\textbackslash ga}), with extra wrapper packages (bm, bbm, dsfont, xspace, cancel, mathtools, enumitem). The delivered numbering is the unit's own. Assets: no retained span contains a figure environment, a graphics-inclusion command, a PDF-page inclusion, a table or a TikZ picture; no image file is copied into this package, the package declares no asset dependency and no image file is redistributed with it. Licence: version arXiv:2011.05237v3 of this article is distributed under the Creative Commons Attribution 4.0 International licence, as recorded in the arXiv licence metadata for this exact version and shown on the abstract page https://arxiv.org/abs/2011.05237v3. A full rights notice is delivered at licenses/arxiv-2011.05237-CC-BY-4.0.txt. Characters: every retained excerpt is pure ASCII, so the delivered engine, XeLaTeX with its default Latin Modern text font, reports no missing character.
\begin{lemma}\label{lem:base-change-lemma}
Suppose $\gL$ is the $L$-ideal associated with the data $(A,T,M,N,[-,-])$. Let $\gL_0$ be the $L$-ideal associated with the data $(A_0,T_0,M_0,N_0,[-,-]_0)$ in one of the two following cases.
\begin{enumerate}[label=(\alph*)]
\item We assume $A_0$ is a flat $A$-algebra, and then define $T_0 = T\otimes_A A_0$ and $M_0 = M\otimes_A A_0$ and $N_0 = N\otimes_A A_0$ and finally $[-,-]_0$ to be the scalar extension of $[-,-]$ along $A \rightarrow A_0$. \label{lemma-part:base-change-lemma-flat}
\item We assume that $T_0 \subseteq T$ is an $A$-subalgebra that is a direct factor of $T$ as an $A$-algebra, { meaning that we can write $T=T_0\times T_1$ as $A$-algebras for a suitable $A$-algebra $T_1$}. Define $A_0 = A$ and $M_0 = M\otimes_T T_0 \subseteq M$ and $N_0 = N\otimes_T T_0 \subseteq N$ and $[-,-]_0$ to be the restriction of $[-,-]$ to $M_0 \otimes_A N_0 \subseteq M\otimes_A N$ 
{(the tensor products $-\otimes_TT_0$ are taken with respect to the canonical projection $T\rightarrow T_0$)}.\label{lemma-part:base-change-lemma-factor}
\end{enumerate}
Then, in either case, there is a canonical isomorphism
\begin{equation*}
\gL\otimes_T T_0 \simeq \gL_0.
\end{equation*}
\end{lemma}

\begin{proof}
It is enough in either case to describe a canonical isomorphism $M_T\otimes_T T_0 \simeq (M_0)_{T_0}$, where $M$ is any $T$-module and $T_0$ and $M_0$ depend on the context \ref{lemma-part:base-change-lemma-flat} or \ref{lemma-part:base-change-lemma-factor}.
\begin{enumerate}[label=(\alph*)]
\item Since annihilators commute with flat { base} change, we have natural isomorphisms
\begin{equation}\label{eqn:flat-natural-isos}
M_T\otimes_T T_0 \simeq M_T\otimes_A A_0 \simeq \left((M_T\otimes_A T) \otimes_A A_0\right)[\ga \otimes_A A_0].
\end{equation}
Now let $\ga_0 = \ker(T_0\otimes_{A_0} T_0 \rightarrow T_0)$. Multiplication commutes with all base change, and kernels commute with flat base change, so there is a canonical isomorphism
\begin{equation*}
\ga \otimes_A A_0 \simeq \ga_0.
\end{equation*}
Thus, the natural isomorphism 
\begin{equation*}
(M\otimes_A T) \otimes_A A_0 \simeq M_0\otimes_{A_0} T_0
\end{equation*}
and \eqref{eqn:flat-natural-isos} induces the claimed isomorphism $M_T\otimes_T T_0 \simeq (M_0)_{T_0}$.
\item 
{ Let $T = T_0 \times T_1$ be as in the statement and let $e_0$ and $e_1$ be the corresponding idempotents so that $e_0+e_1=1_T$ (where $1_T$ is the identity element of $T$), $e_0^2=e_0$, $e_1^2=e_1$, $e_1e_2=0$. 
Write $M_1 = e_1M= M\otimes_T T_1$ (where the tensor product is taken with respect to the canonical projection $T\rightarrow T_1$), and note that $M \simeq M_0 \times M_1$ as $T$-modules, thanks to the equation $e_0+e_1=1_T$. Thus we have an isomorphism of $T\otimes_AT$-modules 
\begin{equation}\label{eqn:base-change-sum}
\begin{split} M\otimes_A T &\simeq ( M \otimes_A T_0)\times (M\otimes_A T_1)\\
&\simeq 
\left(( M_0 \otimes_A T_0)\times (M_1\otimes_A T_0)\right)\times \left((M_0\otimes_A T_1)\times (M_1\otimes_AT_1)\right)
\end{split}
\end{equation}
Let $t\in T$ and write it as $t=t_0+t_1$ with $t_0=te_0$ and $t_1=te_1$. 
Since $e_0e_1=0$, we see easily that the action of $t_1\otimes 1$ is trivial on $M_0\otimes_A T_1$ and $M_0\otimes_AT_0$, while $1\otimes t_1$ acts trivially on $M_0\otimes_AT_0$ and $M_1\otimes T_0$; similarly, the action of $t_0\otimes 1$ is trivial on $M_1\otimes_AT_0$ and $M_1\otimes_AT_1$, while $1\otimes t_0$ acts trivially on $M_1\otimes_AT_1$ and $M_0\otimes_AT_1$.   
Using the decomposition \eqref{eqn:base-change-sum}, write $x\in M\otimes_AT$ as $x=((x_{00},x_{10}),(x_{01},x_{11}))$ with $x_{ij}=(e_i\otimes e_j)x$ for $i,j\in\{0,1\}$. Then we have 
\[\begin{split}
(t\otimes 1-1\otimes t)(x)&=(t\otimes 1-1\otimes t)((x_{00},x_{10}),(x_{01},x_{11}))\\
&=((t_0\otimes1-1\otimes t_0)(x_{00},0),(0,(t_1\otimes 1-1\otimes t_1)x_{11}))
\end{split}\] 
and therefore we find that $M_T$ is given by
\begin{equation*}
M_T \simeq \left((M_0)_{T_0} \times\{0\}\right)\times \left(\{0\}\times (M_1)_{T_1}\right)\simeq (M_0)_{T_0}\times (M_1)_{T_1}.
\end{equation*}
Tracing through these isomorphisms, we see that $M_T\otimes_T T_0 \simeq (M_0)_{T_0}$.}
\end{enumerate}
The proofs are complete.
\end{proof}

\end{document}
